\section{函数空间：从点态观测到一致控制}\label{sec:01-03}

拓扑紧性能够从许多有限条件中产生一个极限点，却不保证这个极限仍是连续函数。若
$f_\alpha(x)$ 对每个 $x$ 都收敛，极限函数甚至可能在每一点都与原函数族的整体形状毫无相似之处。
另一方面，只知道一族简单函数能区分任意两点，也还看不出它们为何能够逼近一个任意的连续函数。
本节处理这两道障碍。首先构造支撑可控的连续函数；随后把等度连续性解释为从有限采样向整个紧集传播
误差的机制；最后证明，函数代数的点分离能力经过紧性放大后会成为一致逼近能力。

\subsection{局部紧空间、紧化与连续函数工具}\label{subsec:1-3-1}

从一个很具体的任务开始。设 $K$ 是空间中的紧集，$U$ 是包含 $K$ 的开集。我们希望找到连续函数
$\varphi$，使它在 $K$ 上等于 $1$，在 $U$ 外等于 $0$，而且非零部分不会悄悄贴到 $U$ 的边界上。
在度量空间里，人们常写下距离函数来完成这件事；一般拓扑空间没有距离，必须找出距离公式真正用到的
分离性质。

\begin{example}[有限维局部紧性与 \(\ell^2\) 反例]\label{ex:1-3-1-1}
在 $\R^n$ 中，每个点 $x$ 都有闭包紧的球邻域。若 $G\subset\R^n$ 开且 $x\in G$，可取 $r>0$
使 $\overline{B(x,r)}\subset G$，因此 $G$ 也有同一性质；离散空间中，单点集就是所需的紧邻域。

这种现象在无限维赋范空间里已经会失败。以 $\ell^2$ 为例，若 $0$ 有一个闭包紧的邻域 $W$，则存在
$r>0$ 使 $B(0,r)\subset W$。闭球 $\overline{B(0,r/2)}$ 是 $\overline W$ 的闭子集，因而紧；
经线性缩放，$\ell^2$ 的闭单位球也应紧。然而标准正交向量满足
\[
  \norm{e_j-e_k}_2=\sqrt2\qquad(j\ne k),
\]
所以 $(e_k)$ 没有 Cauchy 子列，更没有收敛子列。这与紧度量空间的序列紧性矛盾。因此 $\ell^2$
在 $0$ 附近没有紧闭包的邻域。一般赋范空间何时局部紧是另一项结果；这里的计算已经提供了所需的
无限维反例。
\end{example}

\begin{definition}[局部紧 Hausdorff 空间]\label{def:1-3-1-1}
Hausdorff 空间 $X$ 称为\emph{局部紧的}，若每个 $x\in X$ 都有一个邻域，其在 $X$ 中的闭包紧。
记
\[
  A\Subset X
\]
表示 $\overline A$ 紧；写 $A\Subset U$ 时，还要求 $\overline A\subset U$。
\end{definition}

局部紧性允许空间整体向无穷远延伸，却要求每一点附近仍可用紧集控制。把所有逃向无穷远的方向合并成
一个点，会把许多局部构造化为紧空间上的分离问题。

\begin{definition}[一点紧化拓扑]\label{def:1-3-1-2}
设 $X$ 是非紧 Hausdorff 空间，在集合 $X^*=X\cup\{\infty\}$ 上规定：$X$ 中原有的开集仍然开；
含 $\infty$ 的集合 $W$ 为开集，当且仅当 $X^*\setminus W$ 是 $X$ 中的紧集。若这个拓扑是
Hausdorff 的，就称 $X^*$ 为 $X$ 的\emph{一点紧化}。
\end{definition}

定义中的条件确实给出一个拓扑。设一族规定的开集之并为 $O$，且其中某个成员 $W$ 含 $\infty$。令
$K=X^*\setminus W$，则 $K$ 是 $X$ 中的紧集。每个规定开集与 $X$ 的交都在 $X$ 中开：对含
$\infty$ 的成员，这是因为其补集紧且在 Hausdorff 空间 $X$ 中闭。因此 $O\cap X$ 在 $X$ 中开，且
\[
 X^*\setminus O=K\cap\bigl(X\setminus(O\cap X)\bigr)
\]
是 $K$ 的闭子集，因而紧。若没有成员含 $\infty$，该并就是 $X$ 中开集的并。有限交方面，
两个 $\infty$-邻域的交之补集是两个紧集之并；而 $X$ 中开集 $O$ 与
$\{\infty\}\cup(X\setminus K)$ 的交为 $O\setminus K$，它开是因为 Hausdorff 空间中的紧集闭。

\begin{proposition}[一点紧化判据]\label{proposition:1-3-1-one-point-compactification}
设 $X$ 是非紧 Hausdorff 空间，并在 $X^*$ 上采用\cref{def:1-3-1-2}中的拓扑。则 $X^*$ 紧，
$X$ 是 $X^*$ 的开稠密子空间，并且
\[
  X^*\text{ 是 Hausdorff 空间}
  \quad\Longleftrightarrow\quad
  X\text{ 是局部紧空间}.
\]
\end{proposition}

\begin{proof}
给定 $X^*$ 的开覆盖，选取其中一个含 $\infty$ 的成员 $W$。集合
$K=X^*\setminus W$ 是 $X$ 中的紧集，故覆盖中有限多个成员已经覆盖 $K$；再添 $W$ 就得到
$X^*$ 的有限子覆盖。于是 $X^*$ 紧。$X$ 是开子空间，并保持原拓扑。又因 $X$ 非紧，每个形如
$\{\infty\}\cup(X\setminus K)$ 的邻域都与 $X$ 相交，所以 $\infty\in\overline X$，即 $X$ 稠密。

先设 $X$ 局部紧。$X$ 内的两点可由原来的 Hausdorff 性分离。给定 $x\in X$，取开邻域 $V$
使 $\overline V$ 紧。于是
\[
  V
  \quad\text{与}\quad
  \{\infty\}\cup(X\setminus\overline V)
\]
是分别包含 $x$ 与 $\infty$ 的不交开集，故 $X^*$ 为 Hausdorff 空间。

反过来设 $X^*$ 为 Hausdorff 空间。给定 $x\in X$，可取不交开集 $V,W\subset X^*$，分别包含
$x$ 与 $\infty$。按拓扑的定义，存在紧集 $K\subset X$ 使
$\{\infty\}\cup(X\setminus K)\subset W$。因此 $V\cap X\subset K$。紧集 $K$ 在 Hausdorff
空间 $X$ 中闭，故 $\overline{V\cap X}^{\,X}\subset K$，这个闭包是紧集。于是 $x$ 有闭包紧的
邻域 $V\cap X$，从而 $X$ 局部紧。
\end{proof}

\begin{definition}[支撑与三类连续函数]\label{def:1-3-1-3}
标量域 $\mathbb F$ 固定为 $\R$ 或 $\C$。对 $f\in C(X,\mathbb F)$，定义
\[
  \supp f=\overline{\{x\in X:f(x)\ne0\}}.
\]
记 $C_b(X)$ 为所有有界连续标量函数组成的空间，$C_c(X)$ 为其中支撑紧的函数组成的子空间。
函数 $f\in C(X,\mathbb F)$ 称为\emph{在无穷远消失}，若对每个 $\varepsilon>0$，集合
\[
  \{x\in X:|f(x)|\geq\varepsilon\}
\]
紧；所有这类函数组成 $C_0(X)$。三个空间均使用上确界范数
$\norm{f}_\infty=\sup_{x\in X}|f(x)|$；当 $X=\varnothing$ 时，约定唯一函数的上确界范数为 $0$。
上述条件蕴含 $f$ 有界，因而这个范数在 $C_0(X)$ 上有限。
\end{definition}

最后一句值得核对。取 $\varepsilon=1$，集合 $K=\{|f|\ge1\}$ 紧，故 $f|_K$ 有界；在 $X\setminus K$
上已有 $|f|<1$。于是 $f$ 在整个 $X$ 上有界。若 $X$ 非紧且局部紧，则
\[
 f\in C_0(X)
 \quad\Longleftrightarrow\quad
 \widetilde f(x)=
 \begin{cases}
   f(x),&x\in X,\\
   0,&x=\infty
 \end{cases}
 \text{ 属于 }C(X^*).
\]
事实上，$\widetilde f$ 在 $\infty$ 连续恰好表示：对每个 $\varepsilon>0$，存在紧集 $K$，使
$|f|<\varepsilon$ 于 $X\setminus K$。若 $f\in C_0(X)$，可直接取
$K_\varepsilon=\{|f|\ge\varepsilon\}$。反过来，若存在这样的紧集 $K$，则
\[
  \{|f|\ge\varepsilon\}\subset K.
\]
左侧由 $f$ 的连续性在 $X$ 中闭，因而是紧集 $K$ 的闭子集，所以 $f\in C_0(X)$。若 $X$ 紧，则每个
连续函数都在无穷远消失，故
$C_0(X)=C(X)$。

\begin{example}[无穷远处的连续性]\label{ex:1-3-1-2}
在 $\R^n$ 上令
\[
  f(x)=\frac{1}{1+\norm{x}}.
\]
当 $0<\varepsilon\le1$ 时，
$\{f\ge\varepsilon\}=\{\norm{x}\le\varepsilon^{-1}-1\}$ 是闭有界集；当 $\varepsilon>1$ 时该集合为空。
所以 $f\in C_0(\R^n)$，其在一点紧化上满足 $f(\infty)=0$。常数函数 $1$ 则不在
$C_0(\R^n)$ 中，因为 $\{|1|\ge1/2\}=\R^n$ 不紧。这个差别不是关于函数是否有界，而是关于它在
逃离每个紧集后是否一致趋于零。
\end{example}

接下来先在紧 Hausdorff 空间上建立所需的分离工具。证明包含两个层次：拓扑的正规性给出嵌套开集，
二进有理数把这些开集编码成一个连续实函数。

\begin{theorem}[紧 Hausdorff 连续分离工具包]\label{theorem:1-3-1-compact-separation}
每个紧 Hausdorff 空间 $L$ 都是正规的，即任意两个不交闭集都可由不交开集分离，并且有以下结论。
\begin{enumerate}
  \item 若闭集 $A,B\subset L$ 不交，则存在 $u\in C(L,[0,1])$，使
        $u|_A=0$ 且 $u|_B=1$。
  \item 若 $A\subset L$ 闭，$M\ge0$，且 $f\in C(A,[-M,M])$，则存在
        $F\in C(L,[-M,M])$ 使 $F|_A=f$。
\end{enumerate}
\end{theorem}

\begin{proof}
先证正规性。给定 $x\notin B$，考虑所有满足下列条件的开集 $Q$：存在包含 $x$ 的开集 $P$，
使 $P\cap Q=\varnothing$。Hausdorff 性说明这些 $Q$ 覆盖 $B$。闭集 $B$ 是紧集，故其中有限多个
$Q_1,\ldots,Q_m$ 覆盖 $B$；对这有限多个 $Q_j$ 分别取相应的 $P_j$。于是
\[
  P=\bigcap_{j=1}^mP_j,\qquad Q=\bigcup_{j=1}^mQ_j
\]
是不交开集，分别包含 $x$ 与 $B$。现在令 $A,B$ 为不交闭集。对每个 $x\in A$ 作上述点与闭集的
分离。所有能够与某个包含 $B$ 的开集分离的开集 $P$ 构成 $A$ 的开覆盖。由 $A$ 的紧性取其中有限个
$P_1,\ldots,P_r$ 覆盖 $A$，并对这有限多个 $P_i$ 取相应的 $Q_i\supset B$；则
\[
  \bigcup_{i=1}^rP_i
  \quad\text{与}\quad
  \bigcap_{i=1}^rQ_i
\]
给出所需分离。

正规性还有一个反复使用的等价形式：若闭集 $F\subset O$ 且 $O$ 开，则存在开集 $V$ 使
\begin{equation}\label{eq:1-3-normal-shrink}
  F\subset V\subset\overline V\subset O.
\end{equation}
为此分离 $F$ 与闭集 $L\setminus O$，设相应不交开集为 $V$ 与 $W$。由
$\overline V\subset L\setminus W\subset O$ 即得\eqref{eq:1-3-normal-shrink}。

现在证明第一项。令 $D=\{k/2^n:0\le k\le2^n,\ n\in\N\}$。先置
$U_1=L\setminus B$，再由\eqref{eq:1-3-normal-shrink}取 $U_0$，使
$A\subset U_0\subset\overline U_0\subset U_1$。逐层插入奇数分点：若第 $n$ 层的开集已经选定，
就在每一对相邻的 $U_{k/2^n}$ 与 $U_{(k+1)/2^n}$ 之间使用
\eqref{eq:1-3-normal-shrink}，选取 $U_{(2k+1)/2^{n+1}}$，使
\[
 \overline U_{k/2^n}
 \subset U_{(2k+1)/2^{n+1}}
 \subset\overline U_{(2k+1)/2^{n+1}}
 \subset U_{(k+1)/2^n}.
\]
于是对任意 $r,s\in D$ 且 $r<s$，都有 $\overline U_r\subset U_s$。定义
\[
  u(x)=\inf\{r\in D:x\in U_r\},
\]
并在右侧集合为空时令 $u(x)=1$。$A\subset U_0$ 给出 $u|_A=0$，而 $B$ 与每个 $U_r$
都不交，故 $u|_B=1$。

只余连续性。对 $0<a<1$，稠密性与定义给出
\[
  \{u<a\}=\bigcup_{\substack{r\in D\\r<a}}U_r.
\]
另一方面，
\[
  \{u>a\}=\bigcup_{\substack{r\in D\\r>a}}(L\setminus\overline U_r).
\]
为验证第二式的正向包含，若 $u(x)>a$，选 $a<r<s<u(x)$；此时 $x\notin U_s$，而
$\overline U_r\subset U_s$，所以 $x\notin\overline U_r$。反向包含则由
$x\notin\overline U_r$ 推出 $x\notin U_t$ 对每个 $t<r$，因而 $u(x)\ge r>a$。
两个集合都是开集；对阈值 $a\le0$ 或 $a\ge1$，严格次水平集与超水平集由 $0\le u\le1$
直接处理。所以 $u$ 连续。这个二进递归只进行可数多层选择；用后文的记号，它采用
$\textsf{DC}$。

证明第二项。$M=0$ 时取零函数。设 $M>0$，先以 $M$ 缩放，归结为 $|f|\le1$。我们使用如下
一次逼近。若 $r\in C(A)$ 且 $|r|\le q$，令
\[
 A_- =\{r\le-q/3\},\qquad A_+=\{r\ge q/3\}.
\]
这两个闭集不交。由第一项取 $v\in C(L,[0,1])$，使 $v=0$ 于 $A_-$、$v=1$ 于 $A_+$，并令
$g=(q/3)(2v-1)$。于是 $\norm{g}_\infty\le q/3$。在 $A_+$ 上，$g=q/3$；在 $A_-$ 上，
$g=-q/3$；在余下点，$|r|<q/3$ 且 $|g|\le q/3$。逐一比较得到
\begin{equation}\label{eq:1-3-tietze-step}
  |r-g|\le 2q/3\qquad\text{于 }A.
\end{equation}

令 $r_0=f$。假设 $r_n\in C(A)$ 已经定义且
$\norm{r_n}_{C(A)}\le(2/3)^n$。把一次逼近用于 $q=(2/3)^n$，选取
$g_n\in C(L)$，再定义
\[
  r_{n+1}=r_n-g_n|_A
  =f-\sum_{k=0}^{n}(g_k|_A).
\]
由\eqref{eq:1-3-tietze-step}，
$\norm{r_{n+1}}_{C(A)}\le(2/3)^{n+1}$，所以归纳确实可以继续，并得到
\[
 \norm{g_n}_\infty\le\frac13\left(\frac23\right)^n,
 \qquad
 \left|f-\sum_{k=0}^{n}g_k\right|
 \le\left(\frac23\right)^{n+1}\quad\text{于 }A.
\]
级数 $\sum_ng_n$ 在 $L$ 上一致收敛为连续函数 $F$；第一个估计给出
$\norm{F}_\infty\le\sum_{n\ge0}\frac13(2/3)^n=1$，第二个估计给出 $F|_A=f$。
最后乘回 $M$，即得所述延拓。这里逐次选取残差逼近函数也采用可数递归选择，集合论标记同为
$\textsf{DC}$。
\end{proof}

有限离散空间为这个构造提供了直接检验：Urysohn 函数就是闭集的 $0/1$ 分配，Tietze 延拓只是在未指定
点选取区间内的数值。二进开集族的作用，是把同一分配改造成连续函数。需要强调的是，一般 Hausdorff
空间未必正规。一个标准反例是 Sorgenfrey 直线 $S$ 的平方，其中 $S$ 的拓扑基由半开区间
$[a,b)$ 组成。若 $x<y$，则 $[x,y)$ 与 $[y,y+1)$ 是不交基本邻域，故 $S$ 为 Hausdorff 空间，
从而 $S^2$ 也 Hausdorff。每个非空基本区间 $[a,b)$ 都含有有理数
$q$（取 $a\le q<b$），所以 $\Q$ 在 $S$ 中稠密，进而 $\Q^2$ 是 $S^2$ 的可数稠密子集。令
\[
 D=\{(x,-x):x\in\R\}
\]
是基数为连续统的闭离散子空间。若 $(a,b)\notin D$，则 $a+b\ne0$：当 $a+b>0$ 时，任意
$\varepsilon>0$ 给出的矩形 $[a,a+\varepsilon)\times[b,b+\varepsilon)$ 内各点坐标和都为正；当
$a+b<0$ 时取 $0<\varepsilon<-(a+b)/2$，矩形内各点坐标和都小于零。故 $D$ 闭。又若
$(y,-y)\in[x,x+\varepsilon)\times[-x,-x+\varepsilon)$，则同时有 $y\ge x$ 与 $y\le x$，所以
$y=x$；故该矩形与 $D$ 只交于 $(x,-x)$，从而 $D$ 离散。若 $S^2$ 正规，则对
$D$ 的每个子集 $E$，离散性说明 $E$ 与 $D\setminus E$ 都在 $D$ 中闭；再由 $D$ 在 $S^2$ 中闭，
它们也在 $S^2$ 中闭。上面 Urysohn 证明的二进开集构造在正规性建立后只使用正规性，不再使用紧性，
故在这个反设下可照搬，并给出一个在 $E$ 上取 $0$、在 $D\setminus E$ 上取 $1$ 的
连续函数。下面把所需的基数矛盾写成显式映射。记
$\mathscr C=C(S^2,\R)$，并定义
\[
 \Theta:\mathscr C\longrightarrow\mathcal P(D),\qquad
 \Theta(f)=\{d\in D:f(d)<1/2\}.
\]
若 $S^2$ 正规，则对每个 $E\subset D$，上述分离函数 $f$ 满足 $\Theta(f)=E$，所以 $\Theta$ 满射。

另一方面，限制映射
\[
 R:\mathscr C\longrightarrow\R^{\Q^2},\qquad R(f)=f|_{\Q^2}
\]
是单射。事实上，若 $f|_{\Q^2}=g|_{\Q^2}$ 而某点 $x$ 满足 $f(x)\ne g(x)$，则连续函数
$|f-g|$ 的开集
\[
 O=\{z\in S^2:|f(z)-g(z)|>\tfrac12|f(x)-g(x)|\}
\]
是 $x$ 的邻域；它与稠密集 $\Q^2$ 相交，和限制相等矛盾。
还可把 $\R^{\Q^2}$ 单射编码进 $\R$：固定枚举
$\Q^2=\{z_n:n\ge1\}$、$\Q=\{q_m:m\ge1\}$ 以及双射
$\pi:\N^2\to\N$。对数列 $a=(a_n)$ 令
\[
 E_a=\{(n,m):q_m<a_n\},\qquad
 c(a)=\sum_{(n,m)\in E_a}2\,3^{-\pi(n,m)}.
\]
有理数的稠密性说明 $a\ne b$ 时 $E_a\ne E_b$；而三进制数位只取 $0,2$，故上式把不同的
$E_a$ 送到不同实数。更直接地，若 $k$ 是两个编码数位的最小差异位置，则该位贡献
$2\cdot3^{-k}$，其后全部数位贡献的绝对值至多
$\sum_{j>k}2\cdot3^{-j}=3^{-k}$，所以差不为零。因此 $c$ 是单射，进而 $\mathscr C$ 可单射进
$\R$。又 $x\mapsto(x,-x)$ 是 $\R$ 到 $D$ 的双射，故存在单射 $j:\mathscr C\to D$。

常值零函数 $\mathbf0$ 属于 $\mathscr C$。利用它定义满射
\[
 p:D\longrightarrow\mathscr C,\qquad
 p(d)=
 \begin{cases}
   j^{-1}(d),&d\in j(\mathscr C),\\
   \mathbf0,&d\notin j(\mathscr C).
 \end{cases}
\]
于是 $\Theta\circ p:D\to\mathcal P(D)$ 满射。但对任何映射
$h:D\to\mathcal P(D)$，对角集
$H=\{d\in D:d\notin h(d)\}$ 都不可能等于任一 $h(d)$，所以这样的满射不存在。这一矛盾证明
$S^2$ 不正规，也说明不能在任意 Hausdorff 空间中无条件使用连续分离工具。

\begin{proposition}[相对紧邻域引理]\label{proposition:1-3-1-1}
设 $X$ 局部紧 Hausdorff，$K\subset U\subset X$，其中 $K$ 紧、$U$ 开。则存在开集 $V$ 使
\[
  K\subset V\Subset U.
\]
\end{proposition}

\begin{proof}
若 $X$ 非紧，就在紧 Hausdorff 空间 $X^*$ 中工作；若 $X$ 紧，就在 $X$ 中工作。记这个紧空间为
$L$。所有满足 $\overline W^{\,L}\subset U$ 的 $L$-开集 $W$ 覆盖 $K$：对每个 $x\in K$，
把\eqref{eq:1-3-normal-shrink}用于闭集 $\{x\}\subset U$ 即可。紧性给出其中有限个
$W_1,\ldots,W_m$ 覆盖 $K$。令
$V=X\cap\bigcup_{j=1}^mW_j$。它是 $X$ 中的开集，包含 $K$，而
\[
 \overline V^{\,X}\subset\bigcup_{j=1}^m\overline{W_j}^{\,L}\subset U.
\]
右侧有限并是 $L$ 中的紧集且包含于 $X$，所以也是 $X$ 中的紧集。故 $V\Subset U$。
\end{proof}

\begin{theorem}[局部 Urysohn 引理]\label{theorem:1-3-1-2}
设 $X$ 局部紧 Hausdorff，$K\subset U\subset X$，其中 $K$ 紧、$U$ 开。则存在
$\varphi\in C_c(X,[0,1])$，使
\[
  \varphi|_K=1,\qquad \supp\varphi\subset U.
\]
\end{theorem}

\begin{proof}
由\cref{proposition:1-3-1-1}取开集 $V$ 使 $K\subset V\Subset U$。若 $X$ 非紧，在 $X^*$
中对不交闭集 $X^*\setminus V$ 与 $K$（按此顺序）使用
\cref{theorem:1-3-1-compact-separation}；若 $X$ 紧，就在 $X$ 中对 $X\setminus V$ 与 $K$
按同一顺序分离。由该定理中“第一闭集取值 $0$、第二闭集取值 $1$”的约定，所得连续函数
$\varphi$ 在 $K$ 上为 $1$，在 $X\setminus V$ 上为 $0$。
限制到 $X$ 后仍记作 $\varphi$。非零集包含于 $V$，故
$\supp\varphi\subset\overline V\subset U$；又 $\overline V$ 紧，所以 $\varphi\in C_c(X)$。
\end{proof}

\begin{example}[支撑受控的截断]\label{ex:1-3-1-3}
在 $\R$ 上取 $K=[-1,1]$、$U=(-2,2)$，局部 Urysohn 引理给出所需截断。这里还能写出一个公式：
\[
 \varphi(x)=
 \begin{cases}
  1,&|x|\le1,\\
  2-|x|,&1<|x|<2,\\
  0,&|x|\ge2.
 \end{cases}
\]
它连续，$\varphi|_K=1$，且 $\supp\varphi=[-2,2]\subset U$ 这一句其实不成立：端点
$\pm2$ 不属于 $U$。这正暴露了手写公式常见的边界疏忽。把中间断点改为 $3/2$，即令第二行在
$1<|x|<3/2$ 上取 $3-2|x|$、在 $|x|\ge3/2$ 时取零，才得到
$\supp\varphi=[-3/2,3/2]\subset U$。抽象定理把这种严格的支撑余量直接写进结论。
\end{example}

\begin{figure}[htbp]
\centering
\begin{tikzpicture}[scale=0.92]
  \draw[book region,fill=BookAccent!3] (0,0) ellipse (5.2cm and 2.4cm);
  \draw[book region,fill=BookAccent!6] (0,0) ellipse (3.9cm and 1.75cm);
  \draw[book region,fill=BookAccent!10] (0,0) ellipse (2.5cm and 1.05cm);
  \draw[book region,fill=BookWarm!12,draw=BookWarm] (0,0) ellipse (1.15cm and 0.45cm);
  \node[book label] at (0,0) {$K$，$\varphi=1$};
  \node[book label] at (0,0.76) {$V$};
  \node[book label] at (0,1.47) {$\supp\varphi\subset U$};
  \node[book label] at (0,2.14) {$X$};
  \node[book label,anchor=west] at (4.1,1.25) {$\varphi=0$};
  \draw[book accent arrow] (1.2,-0.2) -- (2.4,-0.65)
    node[midway,below,book label] {$K\subset V\Subset U$};
\end{tikzpicture}
\caption{相对紧邻域给截断函数留下严格的支撑余量；局部 Urysohn 引理在这些嵌套集合之间构造截断函数。}
\label{fig:1-3-1-cutoff}
\end{figure}

\begin{theorem}[Tietze 延拓的紧支撑版本]\label{theorem:1-3-1-3}
设 $X$ 局部紧 Hausdorff，$K\subset U\subset X$，其中 $K$ 紧、$U$ 开。若
$f\in C(K,\mathbb F)$，则存在 $\widetilde f\in C_c(X,\mathbb F)$，满足
\[
 \widetilde f|_K=f,\qquad
 \supp\widetilde f\subset U,\qquad
 \norm{\widetilde f}_\infty=\norm f_\infty.
\]
\end{theorem}

\begin{proof}
令 $M=\norm f_\infty$。若 $M=0$，取 $\widetilde f=0$。以下设 $M>0$。由
\cref{proposition:1-3-1-1}取 $K\subset V\Subset U$。在紧 Hausdorff 空间 $\overline V$ 上，
实值情形由\cref{theorem:1-3-1-compact-separation}得到延拓 $F$，满足
$F|_K=f$ 且 $|F|\le M$。

复值情形需要多一步。分别延拓 $\Re f$ 与 $\Im f$，得到连续映射
$G:\overline V\to\R^2\cong\C$。定义闭圆盘上的径向收缩
\[
 R_M(z)=
 \begin{cases}
 z,&|z|\le M,\\
 Mz/|z|,&|z|>M.
 \end{cases}
\]
$R_M$ 连续，且在半径 $M$ 的闭圆盘上为恒等映射。因此 $F=R_M\circ G$ 仍延拓 $f$，并满足
$|F|\le M$。这一步说明，分别延拓实部与虚部后不能直接宣称复模长不增；径向收缩才恢复精确范数界。

再由\cref{theorem:1-3-1-2}取 $\psi\in C_c(X,[0,1])$，使 $\psi=1$ 于 $K$ 且
$\supp\psi\subset V$。在 $X$ 上定义
\[
 \widetilde f(x)=
 \begin{cases}
   \psi(x)F(x),&x\in\overline V,\\
   0,&x\notin\overline V.
 \end{cases}
\]
因为 $\supp\psi\subset V$，若 $x\in\partial V$，则
$x\notin\supp\psi$；闭集 $\supp\psi$ 的补集给出 $x$ 的开邻域 $N_x$，且
\[
  \psi|_{N_x}=0.
\]
故分段函数在 $N_x$ 上恒为零，在每个边界点连续；在 $\overline V$ 的内部及其补集内部，连续性直接来自
各自的定义。
它在 $K$ 上等于 $f$，支撑包含于 $\supp\psi\subset V\subset U$，并且
$\norm{\widetilde f}_\infty\le M$。另一方面，限制到 $K$ 给出
$\norm{\widetilde f}_\infty\ge\norm f_\infty=M$，故等号成立。
\end{proof}

\begin{proposition}[有限单位分解]\label{proposition:1-3-1-finite-partition-of-unity}
设 $X$ 局部紧 Hausdorff，紧集 $K$ 被开集 $U_1,\ldots,U_m$ 覆盖。则存在
$\varphi_j\in C_c(X,[0,1])$，使
\[
  \supp\varphi_j\subset U_j\quad(1\le j\le m),
  \qquad
  \sum_{j=1}^m\varphi_j=1\quad\text{于 }K.
\]
\end{proposition}

\begin{proof}
考虑所有满足
\[
 W\text{ 开},\qquad \overline W\text{ 紧},\qquad
 \overline W\subset U_j\text{ 对某个 }j
\]
的开集 $W$。由\cref{proposition:1-3-1-1}，它们覆盖 $K$。紧性给出有限子覆盖
$W_1,\ldots,W_N$。对每个 $i$ 选一个 $j_i$ 使 $\overline W_i\subset U_{j_i}$；这里只作有限次选择。
由局部 Urysohn 引理，存在 $\psi_i\in C_c(X,[0,1])$，使
$\psi_i=1$ 于 $\overline W_i$ 且 $\supp\psi_i\subset U_{j_i}$。令
\[
  \psi_j=\sum_{\{i:j_i=j\}}\psi_i,
  \qquad s=\sum_{j=1}^m\psi_j.
\]
空和取零。则 $s>0$ 于 $K$，且每个 $\psi_j$ 有紧支撑并满足
$\supp\psi_j\subset U_j$。令 $O=\{s>0\}$。再用局部 Urysohn 引理取
$\chi\in C_c(X,[0,1])$，使 $\chi=1$ 于 $K$ 且 $\supp\chi\subset O$。定义
\[
 \varphi_j(x)=
 \begin{cases}
   \chi(x)\psi_j(x)/s(x),&x\in O,\\
   0,&x\notin O.
 \end{cases}
\]
由于 $\supp\chi$ 是 $O$ 的闭子集，分段定义在 $\partial O$ 附近恒为零，所以 $\varphi_j$ 连续。
在 $O$ 上有 $0\le\psi_j\le s$，所以 $0\le\varphi_j\le\chi\le1$；并有
$\sum_j\varphi_j=\chi=1$ 于 $K$。又
$\supp\varphi_j\subset\supp\psi_j\subset U_j$，结论成立。
\end{proof}

\begin{proposition}[$C_c$ 在 $C_0$ 中稠密]\label{proposition:1-3-1-4}
设 $X$ 局部紧 Hausdorff。则 $C_c(X)$ 在 $C_0(X)$ 中关于上确界范数稠密。
\end{proposition}

\begin{proof}
给定 $f\in C_0(X)$ 与 $\varepsilon>0$，集合
$K=\{x:|f(x)|\ge\varepsilon/2\}$ 紧。对 $K\subset X$ 用局部 Urysohn 引理取
$\varphi\in C_c(X,[0,1])$，使 $\varphi=1$ 于 $K$。于是 $\varphi f\in C_c(X)$。在 $K$ 上
$f-\varphi f=0$；在 $X\setminus K$ 上，
\[
 |f-\varphi f|=|1-\varphi|\,|f|<\varepsilon/2.
\]
所以 $\norm{f-\varphi f}_\infty\le\varepsilon/2<\varepsilon$，这就给出稠密性。
\end{proof}

对任意开集 $U\subset X$，局部 Urysohn 引理还给出一个精确的点态公式：
\[
 \mathbf1_U(x)=
 \sup\{\varphi(x):\varphi\in C_c(X,[0,1]),\ \supp\varphi\subset U\}.
\]
若 $x\notin U$，右侧每一项都为零；若 $x\in U$，对紧集 $\{x\}\subset U$ 应用该引理，右侧就达到
$1$。以后讨论 Radon 正则性时，这个公式会把开集的示性函数转化为支撑受控的连续测试函数。

\begin{proposition}[$C_0(X)$ 的完备性]\label{proposition:1-3-1-c0-complete}
设 $X$ 为拓扑空间，标量域为 $\R$ 或 $\C$。则 $C_0(X)$ 在上确界范数下是 Banach 空间。
\end{proposition}

\begin{proof}
设 $(f_n)$ 是 $C_0(X)$ 中的 Cauchy 列。对每个 $x\in X$，标量列 $(f_n(x))$ 是 Cauchy 列，记其
极限为 $f(x)$。给定 $\varepsilon>0$，取 $N$ 使
$\norm{f_n-f_m}_\infty<\varepsilon$ 对所有 $m,n\ge N$ 成立；固定 $n\ge N$ 并令
$m\to\infty$，得到 $|f_n(x)-f(x)|\le\varepsilon$ 对每个 $x$ 成立。因此 $f_n\to f$ 一致，
$f$ 连续且有界。

还须核对无穷远消失。固定 $\varepsilon>0$，取 $n$ 使
$\norm{f-f_n}_\infty<\varepsilon/2$。于是
\[
 \{x:|f(x)|\ge\varepsilon\}
 \subset \{x:|f_n(x)|\ge\varepsilon/2\}.
\]
右侧因 $f_n\in C_0(X)$ 而紧，左侧又是 $X$ 中的闭集，故它作为右侧的闭子集也紧。于是
$f\in C_0(X)$，并且 $f_n\to f$ 于上确界范数。
\end{proof}

支撑的定义中有一个不能省略的闭包。函数在 $U$ 外为零只推出
$\supp f\subset\overline U$，并不推出 $\supp f\subset U$；
\cref{ex:1-3-1-3}的第一个公式正好展示了差别。另一个边界是完备性：$C_c(X)$ 在上确界范数下通常
不完备。在 $\R$ 上取 $f(x)=(1+|x|)^{-1}$，令 $\chi_n$ 在 $[-n,n]$ 上为 $1$、在
$[-n-1,-n]\cup[n,n+1]$ 上线性降为 $0$、在更外侧为 $0$。则
$f_n=\chi_nf\in C_c(\R)$，且
\[
  \norm{f_n-f}_\infty\le\frac{1}{1+n}\longrightarrow0.
\]
因此 $(f_n)$ 在 $C_c(\R)$ 中是 Cauchy 列，但它唯一可能的上确界范数极限 $f$ 没有紧支撑。
完备的自然空间是 $C_0(X)$。

这里得到的截断函数将用于 Radon 正则性、Riesz--Markov 表示以及 $C_c(X)$ 在 $L^p$ 中的稠密性；
它们共同依赖“紧集内取一、指定开集外取零”的支撑控制。

\begin{exercise}[一点紧化中的收敛]
设 $X$ 非紧且局部紧。证明网 $(x_\alpha)$ 在 $X^*$ 中收敛到 $\infty$，当且仅当对每个紧集
$K\subset X$，它最终落在 $X\setminus K$ 中。
\end{exercise}

\begin{exercise}[重建一点紧化判据]
不引用\cref{proposition:1-3-1-one-point-compactification}的证明，分别从“$X$ 局部紧”和“$X^*$
Hausdorff”出发，构造分离 $x$ 与 $\infty$ 的开集或闭包紧的 $x$-邻域，并证明两个方向互逆。
\end{exercise}

\begin{exercise}[局部紧性的遗传]
证明局部紧 Hausdorff 空间的每个开子空间和闭子空间仍局部紧。再证明 $\Q$ 作为 $\R$ 的子空间
不是局部紧的，由此说明任意子空间不保持局部紧性。
\end{exercise}

\begin{exercise}[有限单位分解的直接核对]
在\cref{proposition:1-3-1-finite-partition-of-unity}的构造中，逐项证明分段定义的 $\varphi_j$ 在
$\partial\{s>0\}$ 连续，并证明其支撑不仅紧，而且包含于指定的 $U_j$。
\end{exercise}

\begin{exercise}[$C_0(X)$ 的完备性]
不引用\cref{proposition:1-3-1-c0-complete}的证明，先证明 $C_b(X)$ 在上确界范数下完备，再把
$C_0(X)$ 识别为 $C_b(X)$ 的闭子空间。
\end{exercise}

局部帽函数和无穷远消失条件给出了可用的函数空间；要在其中识别相对紧集，还必须把逐点有界与邻近点间
的共同振荡控制结合起来。

\subsection{等度连续与 \ArzelaAscoli{} 定理}\label{subsec:1-3-2}

紧性逐点地产生极限并不困难。若对每个 $x\in X$，数集 $\{f(x):f\in\mathcal F\}$ 的闭包紧，
乘积紧性会在 $\prod_{x\in X}\overline{\mathcal F(x)}$ 中给出一个点态聚点。真正的问题是：这个聚点
为何连续，点态收敛又为何能够升级为一致收敛？下面的尖峰把缺少的条件画得很清楚。

Ascoli 的定理是函数空间中对有限维 Heine--Borel 原理的一种早期替代：有限维里“有界”已经很强，
函数空间里还必须控制函数在邻近点之间怎样变化。常见的对角子列证明依靠可分性；下面先用网与乘积紧性
证明一般紧 Hausdorff 版本，再在紧度量空间中给出可计算的有限网格证明。

\begin{example}[移动尖峰]\label{ex:1-3-2-2}
在 $[0,1]$ 上令
\[
  f_n(x)=\max\{1-n|x-1/n|,0\}.
\]
这些函数满足 $0\le f_n\le1$，且 $f_n(1/n)=1$。对 $x=0$，有 $f_n(0)=0$；对固定的 $x>0$，
当 $n>2/x$ 时，
\[
 n|x-1/n|>nx/2>1,
\]
所以 $f_n(x)=0$。因此 $f_n$ 逐点收敛到零函数，但
$\norm{f_n}_\infty=1$，不存在一致收敛到零的子列。事实上它们在 $0$ 处不能共享连续性尺度：
$|1/n-0|\to0$，而 $|f_n(1/n)-f_n(0)|=1$。一致有界没有阻止尖峰越来越窄。
\end{example}

\begin{definition}[等度连续]\label{def:1-3-2-1}
设 $X$ 为拓扑空间，$(Y,d_Y)$ 为度量空间。函数族
$\mathcal F\subset C(X,Y)$ 称为在 $x\in X$ \emph{等度连续}，若对每个
$\varepsilon>0$，存在 $x$ 的邻域 $U$，使
\[
 d_Y(f(y),f(x))<\varepsilon
 \qquad(f\in\mathcal F,\ y\in U).
\]
若它在每个 $x$ 都等度连续，就简称 $\mathcal F$ 等度连续。

当 $X$ 为度量空间且函数为标量值时，记
\[
 \omega_{\mathcal F}(\delta)
 =\sup_{f\in\mathcal F}\sup_{d_X(x,y)\le\delta}|f(x)-f(y)|.
\]
若 $X$ 紧，则等度连续等价于
$\omega_{\mathcal F}(\delta)\to0$（$\delta\downarrow0$）。
\end{definition}

说明最后的等价。共同模量趋于零直接给出各点的等度连续。反过来，给定 $\varepsilon>0$，对每个
$x\in X$ 可取 $r_x>0$，使 $d_X(x,y)<2r_x$ 时所有 $f\in\mathcal F$ 都满足
$|f(y)-f(x)|<\varepsilon/2$。有限多个 $B(x_i,r_{x_i})$ 覆盖 $X$。令
$\delta=\min_i r_{x_i}$。若 $d_X(y,z)\le\delta$，选 $i$ 使 $y\in B(x_i,r_{x_i})$，则
$z\in B(x_i,2r_{x_i})$，从而
\[
 |f(y)-f(z)|
 \le |f(y)-f(x_i)|+|f(x_i)-f(z)|<\varepsilon.
\]

\begin{definition}[点态相对紧]\label{def:1-3-2-2}
若对每个 $x\in X$，集合
$\mathcal F(x)=\{f(x):f\in\mathcal F\}$ 在 $Y$ 中的闭包紧，就称 $\mathcal F$ \emph{点态相对紧}。
标量值函数族点态相对紧，当且仅当它在每一点的取值集合有界。
\end{definition}

\begin{definition}[共同在无穷远消失]\label{def:1-3-2-3}
设 $X$ 局部紧，$\mathcal F\subset C_0(X)$。若对每个 $\varepsilon>0$，存在紧集 $K\subset X$，
使
\[
 |f(x)|<\varepsilon
 \qquad(f\in\mathcal F,\ x\in X\setminus K),
\]
就称 $\mathcal F$ \emph{共同在无穷远消失}。
\end{definition}

若 $X$ 紧而 $Y$ 为度量空间，连续函数 $f,g:X\to Y$ 的距离函数
$x\mapsto d_Y(f(x),g(x))$ 在紧集上有界。因此
\[
  d_\infty(f,g)=\sup_{x\in X}d_Y(f(x),g(x))
\]
是 $C(X,Y)$ 上的度量（$X=\varnothing$ 时仍约定上确界为 $0$）；它诱导的拓扑称为
\emph{一致拓扑}，而 $d_\infty(f_\alpha,f)\to0$ 就是一致收敛。后文紧域上的“一致相对紧”均相对于
这个度量。

\begin{proposition}[等度连续极限升级]\label{proposition:1-3-2-1}
设 $(Y,d_Y)$ 为度量空间，$(f_\alpha)$ 是 $C(X,Y)$ 中的网，并且函数族
$\{f_\alpha\}$ 等度连续。若 $f_\alpha(x)\to f(x)$ 对每个 $x\in X$ 成立，则 $f$ 连续；若 $X$ 紧，
则 $f_\alpha\to f$ 一致收敛。
\end{proposition}

\begin{proof}
固定 $x\in X$ 与 $\varepsilon>0$。等度连续性给出 $x$ 的邻域 $U$，使
\[
 d_Y(f_\alpha(y),f_\alpha(x))<\varepsilon/3
 \qquad(y\in U,\ \alpha).
\]
对固定的 $y\in U$，令 $\alpha$ 趋向有向集的尾部；度量的连续性给出
\[
 d_Y(f(y),f(x))
 =\lim_\alpha d_Y(f_\alpha(y),f_\alpha(x))
 \le\varepsilon/3<\varepsilon.
\]
所以 $f$ 在 $x$ 连续。

再设 $X$ 紧。对每个 $x\in X$，由等度连续性与刚证得的 $f$ 的连续性，存在邻域 $U_x$，使对
$y\in U_x$ 与所有 $\alpha$ 都有
\[
 d_Y(f_\alpha(y),f_\alpha(x))<\varepsilon/3,
 \qquad
 d_Y(f(y),f(x))<\varepsilon/3.
\]
取有限子覆盖 $U_{x_1},\ldots,U_{x_m}$。逐点收敛在这有限多个点上可同时实现：存在 $\alpha_0$，使
$\alpha\ge\alpha_0$ 时
\[
 d_Y(f_\alpha(x_i),f(x_i))<\varepsilon/3
 \qquad(1\le i\le m).
\]
给定 $y\in X$，选 $i$ 使 $y\in U_{x_i}$。三角不等式给出
\[
 d_Y(f_\alpha(y),f(y))
 \le d_Y(f_\alpha(y),f_\alpha(x_i))
   +d_Y(f_\alpha(x_i),f(x_i))
   +d_Y(f(x_i),f(y))
 <\varepsilon.
\]
这个估计对所有 $y$ 同时成立，故收敛一致。
\end{proof}

\begin{theorem}[\ArzelaAscoli{}：紧域]\label{theorem:1-3-2-2}
设 $X$ 为紧 Hausdorff 空间，$(Y,d_Y)$ 为度量空间。若
$\mathcal F\subset C(X,Y)$ 等度连续且点态相对紧，则 $\mathcal F$ 在一致拓扑下相对紧。
特别地，当 $Y$ 紧时，点态相对紧自动成立；标量值情形中，它等价于点态有界。
\end{theorem}

\begin{proof}
对每个 $x\in X$ 令 $K_x=\overline{\mathcal F(x)}\subset Y$。由假设 $K_x$ 紧，故由
Tychonoff 定理，乘积
\[
  E=\prod_{x\in X}K_x
\]
在点态收敛拓扑下紧。把每个 $f\in\mathcal F$ 看作 $E$ 中的点，并令 $H$ 为 $\mathcal F$ 在 $E$
中的闭包；$H$ 是紧集。

先证明 $H$ 仍是等度连续的连续函数族。固定 $h\in H$。由闭包的网刻画，存在
$(f_\alpha)\subset\mathcal F$ 在 $E$ 中收敛到 $h$，也就是逐点收敛到 $h$。由
\cref{proposition:1-3-2-1}，$h$ 连续。再固定 $x\in X$ 与 $\varepsilon>0$。原函数族的等度连续性
给出邻域 $U$，使
\[
 d_Y(f(y),f(x))<\varepsilon/2
 \qquad(f\in\mathcal F,\ y\in U).
\]
对任意 $h\in H$ 取上述逼近网并令其逐点趋于 $h$，便有
\[
 d_Y(h(y),h(x))\le\varepsilon/2<\varepsilon
 \qquad(y\in U).
\]
同一个 $U$ 对所有 $h\in H$ 有效，所以 $H$ 等度连续。

现在考察恒等映射
\[
 \iota:(H,\text{点态拓扑})\longrightarrow(H,\text{一致拓扑}).
\]
若网 $(h_\beta)\subset H$ 在点态拓扑下收敛到 $h\in H$，则刚证得的等度连续性与
\cref{proposition:1-3-2-1}说明它一致收敛到 $h$。由网的连续性判据，$\iota$ 连续。因此点态紧集
$H$ 在一致拓扑下仍紧。

还要确认 $H$ 正是 $\mathcal F$ 的一致闭包。每个 $h\in H$ 都是 $\mathcal F$ 中某网的点态极限，
而该网由\cref{proposition:1-3-2-1}一致收敛到 $h$，故 $H$ 包含于一致闭包。反过来，$H$ 在一致
拓扑下紧，因而在 Hausdorff 的一致拓扑下闭；它包含 $\mathcal F$，所以也包含 $\mathcal F$ 的一致闭包。
两者相等，且这个闭包紧，结论得证。
\end{proof}

乘积紧性只产生坐标极限。上述证明中，等度连续性承担了两项不同工作：先保证极限仍连续，再保证点态拓扑
在闭包上升级为一致拓扑。省略其中任何一步，都会把紧集放在错误的函数空间或错误的拓扑里。

\begin{theorem}[\ArzelaAscoli{}：$C_0$ 版]\label{theorem:1-3-2-3}
设 $X$ 局部紧 Hausdorff。若标量函数族 $\mathcal F\subset C_0(X)$ 等度连续、点态有界，并且共同在
无穷远消失，则 $\mathcal F$ 在上确界范数下相对紧。
\end{theorem}

\begin{proof}
若 $X$ 紧，则 $C_0(X)=C(X)$，结论直接来自紧域版本。以下设 $X$ 非紧。对每个 $f\in\mathcal F$
定义 $\widetilde f\in C(X^*)$，令 $\widetilde f|_X=f$ 且 $\widetilde f(\infty)=0$。
在 $X$ 中各点的等度连续性没有改变；而共同在无穷远消失正好表示：对每个 $\varepsilon>0$，存在
$\infty$ 的邻域 $\{\infty\}\cup(X\setminus K)$，使
\[
 |\widetilde f(x)-\widetilde f(\infty)|<\varepsilon
 \qquad(f\in\mathcal F)
\]
在该邻域成立。因此延拓族在 $X^*$ 上等度连续。它在 $X$ 上点态有界，在 $\infty$ 处恒为零，也就
点态有界。紧域版本说明 $\{\widetilde f:f\in\mathcal F\}$ 在 $C(X^*)$ 中相对紧。延拓保持上确界范数，
而 $\{g\in C(X^*):g(\infty)=0\}$ 是连续评价映射 $g\mapsto g(\infty)$ 的闭核。因此延拓族的闭包
仍在这个闭核中，故 $\mathcal F$ 在 $C_0(X)$ 中相对紧。
\end{proof}

在紧度量空间上，可以把同一个原理写成有限采样算法，不必在证明中显式处理任意乘积。

\begin{proposition}[紧度量空间的有限采样判据]\label{proposition:1-3-2-4}
设 $(X,d)$ 紧，$\mathcal F\subset C(X,\C)$ 点态有界，并且存在函数
$\omega:[0,\infty)\to[0,\infty)$，满足 $\omega(\delta)\to0$（$\delta\downarrow0$）以及
\[
  |f(x)-f(y)|\le\omega(d(x,y))
  \qquad(f\in\mathcal F,\ x,y\in X).
\]
则 $\mathcal F$ 在上确界范数下全有界，因而相对紧。
\end{proposition}

\begin{proof}
若 $X=\varnothing$ 或 $\mathcal F=\varnothing$，结论成立。以下设二者均非空。
给定 $\varepsilon>0$，由 $\omega(t)\to0$ 可选 $\delta>0$，使 $0<t<\delta$ 时都有
$\omega(t)<\varepsilon/3$。紧性给出有限
$\delta$-网 $x_1,\ldots,x_m$。采样映射
\[
  S(f)=(f(x_1),\ldots,f(x_m))\in\C^m
\]
的像有界，因为只涉及有限多个点，而函数族在每一点有界。有限维 Heine--Borel 定理说明
$S(\mathcal F)$ 全有界。以下在 $\C^m$ 上使用最大值范数。先取它的有限 $\varepsilon/6$-网；
对每个与 $S(\mathcal F)$ 相交的
网球，从交集中选一个采样向量，并取一个原像函数 $g_k\in\mathcal F$。这里只有有限多个网球。
三角不等式说明每个 $f\in\mathcal F$ 都有某个 $g_k$ 满足
\[
  \max_{1\le i\le m}|f(x_i)-g_k(x_i)|<\varepsilon/3.
\]
对任意 $x\in X$，取 $i$ 使 $d(x,x_i)<\delta$。若 $x=x_i$，下式首尾两个模量误差均为零；
否则二者都小于 $\varepsilon/3$。于是
\[
 |f(x)-g_k(x)|
 \le |f(x)-f(x_i)|+|f(x_i)-g_k(x_i)|+|g_k(x_i)-g_k(x)|
 <\varepsilon.
\]
所以 $g_1,\ldots,g_N$ 是 $\mathcal F$ 的有限 $\varepsilon$-网。由 Analysis II 中完备度量空间的
全有界判据，$C(X,\C)$ 的完备性把全有界性升级为相对紧性。
\end{proof}

\begin{figure}[htbp]
\centering
\begin{tikzpicture}[node distance=12mm and 13mm]
  \node[book strong node] (grid) {有限 $\delta$-网\\$x_1,\ldots,x_m$};
  \node[book node,right=of grid] (sample) {采样向量\\$S(f)\in\C^m$};
  \node[book node,right=of sample] (reps) {有限代表\\$g_1,\ldots,g_N$};
  \node[book warm node,below=of reps] (uniform) {$\norm{f-g_k}_\infty<\varepsilon$};
  \node[book node,left=22mm of uniform] (modulus) {共同模量\\$\omega(t)<\varepsilon/3\quad(0<t<\delta)$};
  \draw[book arrow] (grid) -- node[above,book label]{观测} (sample);
  \draw[book arrow] (sample) -- node[above,book label]{$\C^m$ 全有界} (reps);
  \draw[book arrow] (reps) -- (uniform);
  \draw[book accent arrow] (modulus) -- node[below,book label]{误差传播} (uniform);
\end{tikzpicture}
\caption{有限采样控制采样向量，共同模量把有限坐标上的误差传播到整个紧集。}
\label{fig:1-3-2-finite-sampling}
\end{figure}

\begin{example}[Lipschitz 球]\label{ex:1-3-2-1}
设 $(X,d)$ 紧，固定 $L,M\ge0$，并令
\[
 \mathcal L_{L,M}
 =\{f\in C(X,\R):|f(x)|\le M,
   \ |f(x)-f(y)|\le Ld(x,y)\text{ 对所有 }x,y\}.
\]
它点态有界，并有共同模量 $\omega(\delta)=L\delta$。由
\cref{proposition:1-3-2-4}，$\mathcal L_{L,M}$ 在 $C(X)$ 中相对紧。它还关于一致收敛闭：若
$f_n\to f$ 一致，则逐点取极限得到 $|f|\le M$ 与
$|f(x)-f(y)|\le Ld(x,y)$。所以 $\mathcal L_{L,M}$ 本身是紧集。
\end{example}

非紧域上还需控制函数向无穷远移动。取
$\phi(x)=\max\{1-|x|,0\}$ 并令 $f_n(x)=\phi(x-3n)$。函数族
$\{f_n\}\subset C_0(\R)$ 共同 $1$-Lipschitz，且在每一点只有至多一个函数非零，因此点态有界；但它不
共同在无穷远消失。不同函数的支撑不交且峰值均为 $1$，故
$\norm{f_n-f_m}_\infty=1$（$n\ne m$）。它没有 Cauchy 子列，从而不相对紧。

点态有界本身也不等于一致有界。令
\[
  c_n=\frac12\left(\frac1{n+1}+\frac1n\right),\qquad
  r_n=\frac12\left(\frac1n-\frac1{n+1}\right),
\]
并在 $[0,1]$ 上定义
\[
  h_n(x)=n\max\left\{1-\frac{|x-c_n|}{r_n},0\right\}.
\]
则 $h_n$ 连续，非零集包含于互不相交的区间 $(1/(n+1),1/n)$，且
\[
  \norm{h_n}_\infty=h_n(c_n)=n.
\]
固定 $x$ 后，至多一个 $h_n(x)$ 非零，所以
$\sup_n|h_n(x)|<\infty$；但上确界范数没有共同界。等度连续性恰好排除了这种越来越陡的峰。

\begin{remark}[积分原函数的共同模量]\label{ex:1-3-2-3}
设 $p>1$、$q=p/(p-1)$，并设函数 $g_n\in L^p([0,1])$ 满足
$\norm{g_n}_{L^p}\le M$。令
\[
  F_n(x)=\int_0^xg_n(t)\,\dd t
\]
。对任意 $x,y\in[0,1]$，记
$I_{x,y}=[\min\{x,y\},\max\{x,y\}]$。\Holder{} 不等式给出
\[
 |F_n(y)-F_n(x)|
 \le \norm{g_n}_{L^p}\norm{\1_{I_{x,y}}}_{L^q}
 \le M|x-y|^{1/q}.
\]
这给出共同 \Holder{} 模量；取 $y=0$ 还得到
$|F_n(x)|\le Mx^{1/q}\le M$，故函数族点态有界。\ArzelaAscoli{} 定理因而给出一致收敛子列。
本评注只说明该紧性定理在积分族上的一种典型应用；$L^p$ 积分与 \Holder{} 不等式见第~3~章。
\end{remark}

同样的紧性结构以后会出现在常微分方程近似解、随机过程的路径以及把有界集合送入相对紧函数族的算子中。
经验过程与 Kolmogorov 连续性定理会再次使用这种机制；Helly 选择原理则依靠单调性或有界变差而不是
等度连续性。上述对象各自还需要新的定义。\ArzelaAscoli{} 只提供函数空间中的紧性判据，并不替代那些
后续论证。

\begin{exercise}[闭包保持等度连续]
设 $X$ 为拓扑空间、$Y$ 为度量空间。证明等度连续族在点态闭包中的所有连续函数仍组成等度连续族；
再说明在紧 $X$ 上，点态闭包与一致闭包在 \ArzelaAscoli{} 的假设下为何相同。
\end{exercise}

\begin{exercise}[一点值控制]
设 $(X,d)$ 紧，函数族共同 $L$-Lipschitz。证明只要在某一点 $x_0$ 的取值集合有界，函数族就在
$X$ 上一致有界，并由此在 $C(X)$ 中相对紧。
\end{exercise}

\begin{exercise}[导数控制的函数族]
证明
\[
 \{f\in C^1([0,1]):|f(0)|\le M,\ \norm{f'}_\infty\le M\}
\]
在 $C([0,1])$ 中相对紧。若删去 $|f(0)|\le M$，说明常数函数为何给出反例。
\end{exercise}

\begin{exercise}[非紧域的缺失条件]
逐项核对平移函数 $f_n(x)=\max\{1-|x-3n|,0\}$ 的等度连续性、点态有界性与
$C_0(\R)$ 成员资格，并直接证明它们不共同在无穷远消失且没有一致收敛子列。
\end{exercise}

有限采样定理告诉我们怎样压缩一个已知函数族。反过来，从少量基本观测函数出发生成整个连续函数空间，
需要把点分离与代数运算转化为一致逼近。

\subsection{Stone--Weierstrass：分离观测生成函数}\label{subsec:1-3-3}

\ArzelaAscoli{} 从有限采样控制一个已经给定的函数族。现在把问题反过来：给定一族基本观测函数，允许作
有限次加法、乘法与极限，它们能否生成所有连续函数？区间上的多项式暗示答案与坐标函数有关，但复数情形
马上提醒我们，单纯“能分离点”还少一个条件。

经典 Weierstrass 定理只处理区间多项式。Stone 的抽象化指出，多项式的具体公式不是核心：含常数、
分离点以及复情形中的共轭封闭，才是把局部插值拼成全局逼近的结构。下面的证明会把这一抽象重新拆回
绝对值、最大值、最小值与两次有限覆盖。

\begin{example}[缺少共轭的失败]\label{ex:1-3-3-3}
令 $\overline{\mathbb D}=\{z\in\C:|z|\le1\}$，并把圆盘代数
$A(\overline{\mathbb D})$ 定义为复多项式在 $C(\overline{\mathbb D},\C)$ 中的一致闭包。它含常数，
坐标函数 $z\mapsto z$ 已经分离闭圆盘中的点。然而 $\overline z$ 不属于这个闭包。

证明只用实变量 Riemann 积分。若多项式 $p_n$ 在闭圆盘上一致逼近 $\overline z$，则在单位圆上
$p_n(e^{it})e^{it}$ 一致趋于 $1$。把复值 Riemann 积分定义为实部与虚部的两个实积分，有
\[
 \left|\int_0^{2\pi}p_n(e^{it})e^{it}\,\dd t-2\pi\right|
 \le 2\pi\sup_{0\le t\le2\pi}
       |p_n(e^{it})-e^{-it}|\longrightarrow0.
\]
另一方面，若 $p(z)=\sum_{k=0}^Na_kz^k$，则
\[
 \int_0^{2\pi}p(e^{it})e^{it}\,\dd t
 =\sum_{k=0}^Na_k\int_0^{2\pi}e^{i(k+1)t}\,\dd t=0,
\]
因为对每个非零整数 $m$，$\cos(mt)$ 与 $\sin(mt)$ 在 $[0,2\pi]$ 上的积分都为零。于是左侧对每个
$n$ 都等于零，不可能趋于 $2\pi$。这个反例没有使用 Cauchy 定理、围道积分或解析函数理论。
\end{example}

\begin{definition}[分离点与处处不消失]\label{def:1-3-3-1}
函数族 $S\subset C(X,\C)$ 称为\emph{分离点}，若对任意 $x\ne y$，存在 $f\in S$ 使
$f(x)\ne f(y)$。若对每个 $x\in X$，存在 $f\in S$ 使 $f(x)\ne0$，就称 $S$ \emph{处处不消失}。
\end{definition}

\begin{definition}[函数代数与 $*$-子代数]\label{def:1-3-3-2}
设 $\mathbb F\in\{\R,\C\}$。函数代数 $A\subset C(X,\mathbb F)$ 是对加法、标量乘法与逐点乘法
封闭的线性子空间；除非另行说明，不要求它含常数。若复函数代数还满足
\[
  f\in A\quad\Longrightarrow\quad\overline f\in A,
\]
就称它为 $*$-子代数。这里 $\overline f(x)=\overline{f(x)}$。实函数代数自动满足对应条件。
\end{definition}

\begin{definition}[一致闭包与函数格]\label{def:1-3-3-3}
函数族 $A\subset C_b(X)$ 关于上确界范数的闭包称为它的\emph{一致闭包}，记作 $\overline A$；
当 $X$ 紧时 $C_b(X)=C(X)$。实函数
线性空间 $L$ 若对逐点的 $\max$ 与 $\min$ 封闭，就称为\emph{函数格}。
\end{definition}

把代数变成函数格是证明的关键。看似需要先知道多项式能逼近绝对值，但若直接调用经典
Weierstrass 定理，就会用待证结论的特例证明一般结论。下面从实幂级数独立构造所需多项式。

\begin{lemma}[绝对值进入闭包]\label{lemma:1-3-3-1}
设 $X$ 紧 Hausdorff，$A\subset C(X,\R)$ 是含常数的函数代数。若 $f\in\overline A$，则
$|f|\in\overline A$。因而 $\overline A$ 是函数格。
\end{lemma}

\begin{proof}
先记 $B=\overline A$。它仍是含常数的函数代数。加法与标量乘法的封闭性来自范数连续性；对于乘法，
若 $f_n\to f$、$g_n\to g$ 一致，则 $(g_n)$ 一致有界，并且
\[
 \norm{f_ng_n-fg}_\infty
 \le\norm{f_n-f}_\infty\norm{g_n}_\infty
   +\norm f_\infty\norm{g_n-g}_\infty\longrightarrow0.
\]

令 $M=\norm f_\infty$。$M=0$ 时结论成立。设 $M>0$。实变量二项级数给出
\begin{equation}\label{eq:1-3-binomial-sqrt}
 \sqrt{1-s}=1-\sum_{n=1}^\infty c_ns^n,
 \qquad 0\le s<1,
\end{equation}
其中
\[
 c_n=\frac{1}{2n-1}\frac{\binom{2n}{n}}{4^n}>0.
\]
还需确认端点 $s=1$ 不破坏一致逼近。对 $0<r<1$，
$\sum_{n\ge1}c_nr^n=1-\sqrt{1-r}\le1$。任意有限部分和令 $r\uparrow1$ 后不超过 $1$，故
$\sum_nc_n\le1$；反过来
$\sum_nc_n\ge\sum_nc_nr^n=1-\sqrt{1-r}$，再令 $r\uparrow1$ 得
$\sum_nc_n\ge1$。所以 $\sum_nc_n=1$。由正项级数的尾和估计，
\eqref{eq:1-3-binomial-sqrt}在 $[0,1]$ 上一致收敛。

定义实多项式
\[
 q_N(t)=M\left[1-\sum_{n=1}^Nc_n
             \left(1-\frac{t^2}{M^2}\right)^n\right].
\]
当 $|t|\le M$ 时，取 $s=1-t^2/M^2\in[0,1]$，由上面的均匀尾和有
\[
 \sup_{|t|\le M}|q_N(t)-|t||
 \le M\sum_{n>N}c_n\longrightarrow0.
\]
因为 $B$ 是含常数代数，$q_N(f)\in B$；一致取极限得到 $|f|\in B$。最后用恒等式
\[
 \max(f,g)=\frac{f+g+|f-g|}{2},
 \qquad
 \min(f,g)=\frac{f+g-|f-g|}{2}
\]
可知 $B$ 对有限的逐点最大值和最小值封闭。
\end{proof}

\begin{theorem}[Stone--Weierstrass：实情形]\label{theorem:1-3-3-2}
设 $X$ 紧 Hausdorff。若 $A\subset C(X,\R)$ 是含常数且分离点的函数代数，则 $A$ 在
$C(X,\R)$ 中一致稠密。
\end{theorem}

\begin{proof}
若 $X=\varnothing$，则 $C(X,\R)$ 只有零函数，结论成立。以下设 $X\ne\varnothing$。
记 $B=\overline A$。由\cref{lemma:1-3-3-1}，$B$ 是函数格。固定目标函数
$f\in C(X,\R)$ 与 $\varepsilon>0$。

对 $x,y\in X$ 构造 $h_{x,y}\in A$，使
\[
 h_{x,y}(x)=f(x),\qquad h_{x,y}(y)=f(y).
\]
当 $x=y$ 时取常数 $f(x)$。当 $x\ne y$ 时，由分离点可取 $a\in A$ 使 $a(x)\ne a(y)$，并置
\[
 h_{x,y}
 =f(x)+\frac{f(y)-f(x)}{a(y)-a(x)}\,[a-a(x)].
\]
因为 $A$ 含常数，$h_{x,y}\in A$，且两个插值等式成立。

先固定 $x$。对每个 $y\in X$，连续函数 $h_{x,y}-f$ 在 $y$ 处为零，所以开集
\[
 V_{x,y}=\{z:h_{x,y}(z)>f(z)-\varepsilon\}
\]
包含 $y$。这些开集覆盖 $X$，由紧性选出 $y_1,\ldots,y_m$ 使
$V_{x,y_1},\ldots,V_{x,y_m}$ 覆盖 $X$。令
\[
 g_x=\max_{1\le j\le m}h_{x,y_j}\in B.
\]
对每个 $z\in X$，至少一个 $V_{x,y_j}$ 包含 $z$，故
$g_x(z)>f(z)-\varepsilon$。另一方面，所有 $h_{x,y_j}(x)$ 都等于 $f(x)$，所以 $g_x(x)=f(x)$。
于是
\[
 U_x=\{z:g_x(z)<f(z)+\varepsilon\}
\]
是包含 $x$ 的开集。

当 $x$ 变化时，$U_x$ 覆盖 $X$。再由紧性选取 $x_1,\ldots,x_r$ 使这些开集覆盖 $X$，并令
\[
  g=\min_{1\le i\le r}g_{x_i}\in B.
\]
每个 $g_{x_i}$ 在全空间都大于 $f-\varepsilon$，故
$g>f-\varepsilon$。给定 $z\in X$，选 $i$ 使 $z\in U_{x_i}$，则
$g(z)\le g_{x_i}(z)<f(z)+\varepsilon$。因此
\[
  |g(z)-f(z)|<\varepsilon\qquad(z\in X).
\]
这说明 $f$ 属于 $B$。目标函数任意，故 $B=C(X,\R)$。
\end{proof}

证明中的两次有限覆盖承担不同方向的不等式。第一次对固定 $x$ 取有限最大值，使函数在全空间位于
$f-\varepsilon$ 上方；第二次取有限最小值，使每一点至少受到一个 $f+\varepsilon$ 上界的控制。
若交换最大值与最小值，两个方向不会自动保留。

\begin{theorem}[Stone--Weierstrass：复情形]\label{theorem:1-3-3-3}
设 $X$ 紧 Hausdorff。若 $A\subset C(X,\C)$ 是含常数、分离点的 $*$-子代数，则 $A$ 在
$C(X,\C)$ 中一致稠密。
\end{theorem}

\begin{proof}
令
\[
 A_{\R}=\{a\in A:a=\overline a\}.
\]
它是 $C(X,\R)$ 中含常数的实函数代数。若 $x\ne y$，取 $a\in A$ 使 $a(x)\ne a(y)$。两个复数
$a(x),a(y)$ 的实部或虚部至少有一个不同，而
\[
 \Re a=\frac{a+\overline a}{2}\in A_{\R},
 \qquad
 \Im a=\frac{a-\overline a}{2i}\in A_{\R}.
\]
所以 $A_{\R}$ 分离点。由实情形，$A_{\R}$ 在 $C(X,\R)$ 中稠密。

给定 $f=u+iv\in C(X,\C)$ 与 $\varepsilon>0$，可取 $a,b\in A_{\R}$ 使
$\norm{u-a}_\infty<\varepsilon/2$、$\norm{v-b}_\infty<\varepsilon/2$。于是 $a+ib\in A$，且
\[
 \norm{f-(a+ib)}_\infty
 \le\norm{u-a}_\infty+\norm{v-b}_\infty<\varepsilon.
\]
故 $A$ 一致稠密。
\end{proof}

\begin{figure}[htbp]
\centering
\begin{tikzpicture}[node distance=11mm and 13mm]
  \node[book strong node] (sep) {分离 $x,y$};
  \node[book node,right=of sep] (interp) {两点插值\\$h_{x,y}$};
  \node[book node,right=of interp] (max) {有限 $\max$\\固定 $x$};
  \node[book node,below=of max] (min) {有限 $\min$\\覆盖 $X$};
  \node[book warm node,left=of min] (approx) {$\norm{g-f}_\infty<\varepsilon$};
  \draw[book arrow] (sep) -- (interp);
  \draw[book arrow] (interp) -- node[above,book label]{格运算} (max);
  \draw[book arrow] (max) -- node[right,book label]{紧性} (min);
  \draw[book accent arrow] (min) -- (approx);
\end{tikzpicture}
\caption{Stone--Weierstrass 证明的两次拼接。最大值与最小值分别保留一侧误差。}
\label{fig:1-3-3-stone-weierstrass}
\end{figure}

\begin{example}[区间多项式]\label{ex:1-3-3-1}
实多项式在 $[a,b]$ 上构成含常数的函数代数。若 $x\ne y$，坐标多项式 $p(t)=t$ 满足
$p(x)\ne p(y)$，所以它分离点。由\cref{theorem:1-3-3-2}，对每个
$f\in C([a,b],\R)$ 与 $\varepsilon>0$，存在实多项式 $p$ 使
$\norm{f-p}_\infty<\varepsilon$。经典 Weierstrass 逼近定理由此成为一般定理的直接实例。
\end{example}

\begin{example}[三角多项式]\label{ex:1-3-3-2}
在 $\mathbb T=\{z\in\C:|z|=1\}$ 上，坐标函数 $\zeta(z)=z$ 与其共轭满足
$\overline\zeta(z)=z^{-1}$。它们生成的 $*$-代数恰为有限 Laurent 和
\[
  \sum_{k=-N}^Na_kz^k,
\]
也就是三角多项式。该代数含常数，且 $\zeta$ 本身分离 $\mathbb T$ 的点。由
\cref{theorem:1-3-3-3}，三角多项式在 $C(\mathbb T,\C)$ 中一致稠密。这里的结论仍然只涉及
连续函数与代数运算；Fourier 系数和复积分尚未出现。
\end{example}

圆盘反例现在也可准确定位：复多项式闭包含常数并分离点，却不对共轭封闭；正是缺少的 $*$ 条件使
$\overline z$ 无法进入闭包。

\begin{theorem}[Stone--Weierstrass：局部紧版本]\label{theorem:1-3-3-4}
设 $X$ 局部紧 Hausdorff，$A\subset C_0(X,\C)$ 是分离点、处处不消失的 $*$-子代数。则 $A$ 在
$C_0(X,\C)$ 中一致稠密。
\end{theorem}

\begin{proof}
先设 $X$ 非紧。把每个 $a\in A$ 以 $a(\infty)=0$ 延拓到 $X^*$，所得代数记为 $\widetilde A$，并令
\[
  A^+=\widetilde A+\C\mathbf1\subset C(X^*,\C).
\]
这是含常数的 $*$-子代数。$A$ 分离 $X$ 内的两点；若 $x\in X$，处处不消失给出 $a\in A$ 使
$a(x)\ne0=a(\infty)$，所以 $A^+$ 也分离 $x$ 与 $\infty$。紧空间复版本说明 $A^+$ 在
$C(X^*)$ 中稠密。

给定 $f\in C_0(X)$，把它延拓为 $\widetilde f\in C(X^*)$，其中 $\widetilde f(\infty)=0$。给定
$\varepsilon>0$，取 $a\in A$ 与 $c\in\C$，使
\[
 \norm{\widetilde f-(\widetilde a+c\mathbf1)}_\infty<\varepsilon/2.
\]
在 $\infty$ 处代入得到 $|c|<\varepsilon/2$，故
\[
 \norm{f-a}_\infty
 \le\norm{\widetilde f-(\widetilde a+c\mathbf1)}_\infty+|c|
 <\varepsilon.
\]

若 $X$ 本身紧，就令 $Y=X\sqcup\{\infty\}$，并把新点规定为孤立点。$Y$ 是紧 Hausdorff 空间，
但这里不把它称为 $X$ 的一点紧化，因为 $X$ 在 $Y$ 中不稠密。把 $A$ 中函数在新点取零，并把所得代数
记为 $\widetilde A$。代数 $\widetilde A+\C\mathbf1$ 用 $A$ 分离 $X$ 中的两点；给定 $x\in X$，
处处不消失性给出 $a\in A$ 使 $a(x)\ne0=\widetilde a(\infty)$，所以它也分离 $x$ 与 $\infty$。
紧空间复版本表明该含常数 $*$-代数在 $C(Y)$ 中稠密。把 $f\in C(X)=C_0(X)$ 在 $\infty$ 处延拓为
零，并取
\(\norm{\widetilde f-(\widetilde a+c\mathbf1)}_\infty<\varepsilon/2\)。在 $\infty$ 处代入得
$|c|<\varepsilon/2$，限制回 $X$ 后便有
$\norm{f-a}_\infty<\varepsilon$，结论成立。
\end{proof}

局部紧版本不要求 $A$ 含常数。在非紧空间中，非零常数通常不属于 $C_0(X)$；处处不消失取代了含常数
所承担的一部分作用，因为它恰好把空间中的每个点与 $\infty$ 分开。分离点也只是全局的观测条件，
与导数是否非零、是否提供局部坐标没有等价关系。

\begin{corollary}[紧度量空间上的可数稠密测试族]\label{corollary:1-3-3-separable}
若 $(X,d)$ 是紧度量空间，则 $C(X,\R)$ 与 $C(X,\C)$ 都可分。
\end{corollary}

\begin{proof}
对每个 $n$ 取有限 $1/n$-网，所有这些网的并给出可数稠密集 $D\subset X$；按后文的记号，这个
逐个选取有限网的步骤使用 $\textsf{CC}$。对 $q\in D$ 定义
$h_q(x)=d(x,q)$。这族函数分离点：若 $x\ne y$，令 $\rho=d(x,y)>0$，取
$q\in D$ 使 $d(x,q)<\rho/3$，则
\[
 |h_q(y)-h_q(x)|
 \ge d(x,y)-2d(x,q)>\rho/3.
\]
由 $1$ 与所有 $h_q$ 生成的实函数代数依
\cref{theorem:1-3-3-2}在 $C(X,\R)$ 中稠密。只允许有理系数、有限多个生成元的多项式仍构成可数集。
具体地，固定
\[
  P=\sum_{\alpha\in F}c_\alpha
       h_{q_1}^{\alpha_1}\cdots h_{q_m}^{\alpha_m},
\]
其中 $F\subset\N^m$ 有限，并给定 $\varepsilon>0$。令
$M_\alpha=\prod_{j=1}^m\norm{h_{q_j}}_\infty^{\alpha_j}$（约定 $0^0=1$），并选
$r_\alpha\in\Q$ 使
\[
  \sum_{\alpha\in F}|c_\alpha-r_\alpha|M_\alpha<\varepsilon.
\]
则相应有理系数多项式 $Q$ 满足
\[
  \norm{P-Q}_\infty
  \le\sum_{\alpha\in F}|c_\alpha-r_\alpha|M_\alpha<\varepsilon.
\]
故这个可数子集已经在 $C(X,\R)$ 中稠密。复情形改用有理实部与有理虚部的系数即可。
\end{proof}

以下只作后文术语预告；Radon 测度与积分将在测度论章节定义，本章不以这段话为任何证明的前提。
对两个有限 Radon 测度 $\mu,\nu$，若它们在一个含常数、分离点的 $*$-代数 $A$ 上给出相同积分，
则 Stone--Weierstrass 定理允许对任意连续函数 $f$ 取 $g_n\in A$ 使
$\norm{g_n-f}_\infty\to0$。有限性给出显式估计
\[
 \left|\int f\,\mathrm d\mu-\int f\,\mathrm d\nu\right|
 \le \mu(X)\norm{f-g_n}_\infty
     +\nu(X)\norm{f-g_n}_\infty
     +\left|\int g_n\,\mathrm d\mu-\int g_n\,\mathrm d\nu\right|.
\]
最后一项为零，令 $n\to\infty$ 即把积分相等性推广到全部连续函数。这一识别方法将在相应概念建立后
用于矩问题与 Fourier 分析。

\begin{exercise}[稠密性的必要条件]
设 $X$ 局部紧 Hausdorff。证明若 $A\subset C_0(X)$ 一致稠密，则 $A$ 分离点且处处不消失。
反过来，指出还必须给 $A$ 增加哪些代数封闭条件，才能应用局部紧 Stone--Weierstrass 定理。
\end{exercise}

\begin{exercise}[三角多项式]
不省略代数验证，证明所有有限和 $\sum_{k=-N}^Na_kz^k$ 在 $\mathbb T$ 上构成含常数、分离点的
$*$-子代数，并由复 Stone--Weierstrass 定理推出其一致稠密性。
\end{exercise}

\begin{exercise}[多变量多项式]
设 $K=\prod_{j=1}^n[a_j,b_j]\subset\R^n$。证明实多变量多项式限制到 $K$ 后构成含常数且分离点的
函数代数，从而在 $C(K,\R)$ 中一致稠密。
\end{exercise}

\begin{exercise}[可数分离族与可度量化]
证明紧 Hausdorff 空间 $X$ 可度量，当且仅当存在可数族 $\{f_n\}\subset C(X,\R)$ 分离点。
提示：一个方向使用距离函数；另一个方向把 $X$ 映入可数乘积
$\prod_n[-\norm{f_n}_\infty,\norm{f_n}_\infty]$，并证明连续单射在紧空间到 Hausdorff 空间之间
给出到其像的同胚。
\end{exercise}

本章还需回答一个隐藏在上述证明里的问题。乘积紧性使用怎样的选择原则？何时可用对角序列代替网？
完备性、可分性与可数拓扑基将决定这些替代何时成立。下一节就整理这些可数性边界。
